\paragraph*{Rarefaction waves}

We now construct a family of continuous waves that generalize the rarefaction waves for the single conservation law. As in the case of shocks, we prove the existence of $n$ curves emanating from a left state $\bm{u}^l$ giving the possible right states $\bm{u}^r$ that can be connected directly using a single rarefaction wave.

The general idea is based on the construction for the single conservation law. Suppose we have a situation where for some $k=1,2,\ldots,n$ we have
\begin{equation}\label{eq:rarefaction-increasing-characteristic-speed}
\lambda_k(\bm{u}^l)<\lambda_k(\bm{u}^r).
\end{equation}
Note that the Lax condition immediately rules out connecting the two states with a single shock. We now mimic the procedure followed for the single equation case and draw characteristic lines $x=\lambda(\bm{u}^l)t$ and $x=\lambda(\bm{u}^r)t$ emanating from the left and right of the origin. (Note that in the case of systems it is not necessary that solutions be constant along characteristics, though in this case, such a "guess" will lead us to a solution.) Observe that this characteristic diagram is very similar to Figure~\ref{fig:characteristic-gap}: we have two regions in the upper half of the $(x,t)$-plane covered by characteristics with a wedge-shaped blank region in between. If we yield to temptation and define a solution to be the constant $\bm{u}^l$ in the left-hand shaded region and $\bm{u}^r$ in the right-hand region, how are we to fill in the blank region? The answer is that we can do so with the following type of wave.

\begin{Definition}
\label{def:system-simple-wave}
Let $\bm{u}$ be a $C^1$ solution of conservation law \eqref{eq:conservation-law-system} in a domain $D$. Then $\bm{u}$ is said to be a \textbf{k-rarefaction wave} (or a \textbf{k-simple wave}) if all $k$-Riemann invariants are constant in $D$.
\end{Definition}

As we might have hoped from observing the results for the single conservation law, if we can find a rarefaction wave that fills in the blank wedge, the characteristics associated with $\lambda_k$ form a "fan."

\begin{Theorem}
\label{thm:system-rarefaction-straight-characteristics}
Let $\bm{u}$ be a $k$-rarefaction wave in a domain $D$. Then the characteristic curves $\hat x'(t)=\lambda_k(\bm{u}(\hat x(t),t))$ are straight lines along which $\bm{u}$ is constant.
\end{Theorem}

\begin{Proof}
We wish to show
\begin{equation}\label{eq:rarefaction-constant-along-characteristics}
0=\dfrac{\mathrm{d}}{\mathrm{d}t}\bm{u}(\hat x(t),t)=\bm{u}_t+\lambda_k\bm{u}_x.
\end{equation}
Lemma~\ref{lem:riemann-invariants-existence} asserts that there exist $n-1$ $k$-Riemann invariants $w_i$, $i=1,2,\ldots,n-1$, whose gradients are linearly independent. Since $\bm{u}$ is a $k$-rarefaction wave, $w_i(\bm{u}(x,t))$ is constant, and hence
\begin{equation}\label{eq:rarefaction-invariant-directional-derivatives}
\dfrac{\mathrm{d}}{\mathrm{d}t}w_i(\bm{u}(\hat x(t),t))=\nabla w_i\cdot(\bm{u}_t+\lambda_k\bm{u}_x)=0,
\end{equation}
for $i=1,2,\ldots,n-1$. We now use the fact that $\bm{u}$ solves \eqref{eq:conservation-law-system} and the definition of $\bm{l}_k$ to deduce
\begin{equation}\label{eq:rarefaction-left-eigenvector-identity}
\bm{l}_k\cdot(\bm{u}_t+\lambda_k\bm{u}_x)=\bm{l}_k\cdot(\bm{u}_t+\nabla\bm{f}\bm{u}_x)=0.
\end{equation}
Thus, $\bm{u}_t+\lambda_k\bm{u}_x$ is orthogonal to every vector in the set
\begin{equation}\label{eq:rarefaction-invariant-gradient-basis}
V:=\{\bm{l}_k,\nabla w_1,\nabla w_2,\ldots,\nabla w_{n-1}\}.
\end{equation}
Thus, all that remains to complete the proof is to show that $V$ is a basis for $\mathbb{R}^n$. This is left to the reader (cf. Problem~\ref{ex:riemann-invariant-gradient-basis}).
\end{Proof}

We now state our basic theorem on the existence of rarefaction curves.

\begin{Theorem}
\label{thm:rarefaction-curves-existence}
Suppose that the system of conservation laws \eqref{eq:conservation-law-system} is genuinely nonlinear in an open region $\Omega\subset\mathbb{R}^n$ in state space, and let the right eigenvectors $\bm{r}_k$ be normalized so that $\nabla\lambda_k\cdot\bm{r}_k=1$. Then for any left state $\bm{u}^l\in\Omega$ there exist $n$ intervals $J_k=[0,a_k)$ and $n$ smooth, one-parameter families of right states $\hat{\bm{u}}_k(\gamma)$ defined for $\gamma\in J_k$ that can be connected to $\bm{u}^l$ by a $k$-simple wave using the procedure above. Moreover, these one-parameter families satisfy the following properties:
\begin{equation}\label{eq:rarefaction-curve-base-state}
\tilde{\bm{u}}_k(0)=\bm{u}^l,
\end{equation}

\begin{equation}\label{eq:rarefaction-curve-initial-tangent}
\tilde{\bm{u}}_k'(0)=\bm{r}_k,
\end{equation}
and for $0<\gamma\in J_k$,
\begin{equation}\label{eq:rarefaction-curve-increasing-speed}
\lambda_k(\bm{u}^l)<\lambda_k(\tilde{\bm{u}}_k(\gamma)).
\end{equation}

\end{Theorem}

\begin{Proof}
The rarefaction curves are simply solutions of the ODE initial-value problem
\begin{equation}\label{eq:rarefaction-curve-ode}
\dfrac{\mathrm{d}\tilde{\bm{u}}_k}{\mathrm{d}\gamma}(\gamma)=\bm{r}_k(\tilde{\bm{u}}_k(\gamma)),
\end{equation}

\begin{equation}\label{eq:rarefaction-curve-ode-initial-data}
\tilde{\bm{u}}_k(0)=\bm{u}^l.
\end{equation}
Existence on an interval about $0$ follows from Theorem~\ref{thm:ode-existence}. Note that
\begin{equation}\label{eq:rarefaction-normalized-speed-derivative}
\dfrac{\mathrm{d}}{\mathrm{d}\gamma}\lambda_k(\tilde{\bm{u}}_k(\gamma))=\nabla\lambda_k\cdot\tilde{\bm{u}}_k'=\nabla\lambda_k\cdot\bm{r}_k=1.
\end{equation}
Thus, $\gamma\mapsto\lambda_k(\tilde{\bm{u}}_k(\gamma))$ is increasing so that \eqref{eq:rarefaction-curve-increasing-speed} holds. Moreover, using the initial condition \eqref{eq:rarefaction-curve-base-state}, we get
\begin{equation}\label{eq:rarefaction-normalized-speed}
\lambda_k(\tilde{\bm{u}}_k(\gamma))=\gamma+\lambda_k(\bm{u}^l).
\end{equation}
To see that we can use this curve to "hook up" a left and right state using a $k$-rarefaction wave, we simply let
\begin{equation}\label{eq:rarefaction-similarity-wave}
\bm{u}(x,t):=\tilde{\bm{u}}_k(x/t-\lambda_k(\bm{u}^l)).
\end{equation}
Note that this is indeed a solution of \eqref{eq:conservation-law-system} and that for any $k$-Riemann invariant $w_k$
\begin{equation}\label{eq:rarefaction-source-time-derivative}
\dfrac{\partial}{\partial t}w_k(\bm{u}(x,t))=\nabla w_k\cdot\tilde{\bm{u}}_k'\left(\dfrac{x}{t^2}\right)=\nabla w_k\cdot\bm{r}_k\left(\dfrac{x}{t^2}\right)=0.
\end{equation}
A similar calculation for the derivative with respect to $x$ shows that any $k$-Riemann invariant is constant in the "fan" region so that the solution is a $k$-rarefaction wave.
\end{Proof}

\begin{figure}[H]
\centering
\includegraphics[width=.78\linewidth]{figures/p-system-rarefaction-curves-source.png}
\caption{}\label{fig:p-system-rarefaction-curves}
\end{figure}

In the case of the $p$ system it is easier to solve \eqref{eq:rarefaction-curve-ode} without normalizing the eigenvectors. We get the curves
\begin{equation}\label{eq:p-system-first-rarefaction-curve}
R_1:\quad\tilde{\bm{u}}_1(\gamma)=\begin{pmatrix}\tilde w_1(\gamma)\\\tilde v_1(\gamma)\end{pmatrix}:=\begin{pmatrix}w^l+\gamma\\v^l+\int_0^\gamma\sqrt{-p'(w^l+\xi)}\,\mathrm{d}\xi\end{pmatrix}
\end{equation}

\begin{equation}\label{eq:p-system-second-rarefaction-curve}
R_2:\quad\tilde{\bm{u}}_2(\gamma)=\begin{pmatrix}\tilde w_2(\gamma)\\\tilde v_2(\gamma)\end{pmatrix}:=\begin{pmatrix}w^l-\gamma\\v^l+\int_0^\gamma\sqrt{-p'(w^l-\xi)}\,\mathrm{d}\xi\end{pmatrix}.
\end{equation}


Because we have not normalized the eigenvectors it is somewhat harder to determine $(w(x,t),v(x,t))$ from the rarefaction curves. To compute a 1-wave between $\bm{u}^l$ and $\bm{u}^r$ (with $\bm{u}^r$ on $R_1$) we take
\begin{equation}\label{eq:p-system-first-rarefaction-speed-range}
\lambda_1(\bm{u}^l)<\dfrac{x}{t}<\lambda_1(\bm{u}^r),
\end{equation}
and solve the equation
\begin{equation}\label{eq:p-system-first-rarefaction-volume}
\dfrac{x}{t}=\lambda_1=-\sqrt{-p'(w(x/t))}
\end{equation}
for $w(x/t)$. Next, we let
\begin{equation}\label{eq:p-system-first-rarefaction-parameter}
\gamma=w(x/t)-w^l,
\end{equation}
and use this in \eqref{eq:p-system-first-rarefaction-curve} to determine $v(x/t)$.

The picture here is much the same as the one for the shock curves. The two rarefaction curves emanate from the left state $\bm{u}^l$; they share the tangent vectors $\bm{r}_k$ with the shock curves, but propagate in the opposite direction. The $R_1$ curve represents rarefaction waves in which both the left and right state have negative speed (sometimes called a slow wave or a back-wave) whereas the $R_2$ curve represents rarefaction waves in which both the left and right state have positive speed (fast waves or front-waves).

\begin{figure}[H]
\centering
\includegraphics[width=.78\linewidth]{figures/slow-fast-state-curves-source.png}
\caption{The slow curve $(S_1\cup R_1)$ and the fast curve $(S_2\cup R_2)$.}\label{fig:slow-fast-state-curves}
\end{figure}

\paragraph*{General solution}

We now show how shocks and rarefaction waves can be put together to get a general solution for the Riemann problem. Our basic theorem is as follows.

\begin{Theorem}
\label{thm:local-riemann-solution}
Suppose that our system of conservation laws \eqref{eq:conservation-law-system} is strictly hyperbolic and genuinely nonlinear in a region $\Omega\subset\mathbb{R}^n$ of state space. Then for any $\bm{u}^l\in\Omega$ there is a neighborhood $N\subset\Omega$ of $\bm{u}^l$ such that if $\bm{u}^r\in N$ there exists a weak solution of the Riemann problem \eqref{eq:conservation-law-system}, \eqref{eq:riemann-piecewise-constant-data}. This solution is composed of at most $n+1$ constant states separated by rarefaction waves and shocks satisfying the Lax shock condition.
\end{Theorem}

We will not prove this, but we show how the process works in the case of the $p$ system. We start with the left state $\bm{u}^l$ and consider the two pairs of shock and rarefaction curves emanating from that point. It is best to think of these as being two $C^1$ curves: a "slow curve" consisting of the union of $S_1$ and $R_1$ (the slow shocks and the slow rarefaction waves) and a "fast curve" consisting of the union of $S_2$ and $R_2$. Of course, if the right state $\bm{u}^r$ in our Riemann problem lies on either of these curves, we can simply connect the left and right state with a single wave (slow or fast, shock or rarefaction), depending on which of the four original curves it lies on. The question remains, what happens if the right state lies in one of the four open regions cut out by our curves?

The solution is obtained by covering these four regions with a family of fast curves. Through each point $\bar{\bm{u}}$ on the slow curve through $\bm{u}^l$ we can construct the curves $\bar S_2$ and $\bar R_2$. These new curves will represent shocks and rarefaction waves, respectively, all with positive speed and all having $\bar{\bm{u}}$ as the left state. Taking the union of $\bar S_2$ and $\bar R_2$ gives us a family of curves (which we will call the "fast family") parameterized by the points $\bar{\bm{u}}$ on the original slow curve. It is left as an exercise (cf. Problem~\ref{ex:p-system-fast-family-covering}) to show that there is a neighborhood $\mathcal N$ of $\bm{u}^l$ that is covered univalently by the fast family; i.e., for any point $\bm{u}^r\in\mathcal N$ there is exactly one member of the fast family containing $\bm{u}^r$.

\begin{figure}[H]
\centering
\includegraphics[width=.78\linewidth]{figures/slow-curve-fast-family-source.png}
\caption{Slow curve and fast family.}\label{fig:slow-curve-fast-family}
\end{figure}

\begin{figure}[H]
\centering
\includegraphics[width=.78\linewidth]{figures/slow-rarefaction-fast-shock-source.png}
\caption{Slow rarefaction--fast shock.}\label{fig:slow-rarefaction-fast-shock}
\end{figure}

\begin{figure}[H]
\centering
\includegraphics[width=.78\linewidth]{figures/slow-shock-fast-shock-source.png}
\caption{Slow shock--fast shock.}\label{fig:slow-shock-fast-shock}
\end{figure}

Now that we have used the left state to generate the slow curve and the fast family, the solution of the Riemann problem is simple. From any right state $\bm{u}^r\in\mathcal N$ we simply follow the appropriate member of the fast family back to a point $\bar{\bm{u}}$ on the slow curve. The point $\bar{\bm{u}}$ is now used as an intermediate state between two waves: a slow wave connecting $\bm{u}^l$ and $\bar{\bm{u}}$ and a fast wave connecting $\bar{\bm{u}}$ and $\bm{u}^r$. Of course, each of the two waves can be either a shock or rarefaction wave depending on which of the four regions defined by the original slow and fast curves the right state $\bm{u}^r$ lies in. The various possibilities are described in Figures~\ref{fig:slow-rarefaction-fast-shock}--\ref{fig:slow-rarefaction-fast-rarefaction}.

\begin{figure}[H]
\centering
\includegraphics[width=.78\linewidth]{figures/slow-shock-fast-rarefaction-source.png}
\caption{Slow shock--fast rarefaction.}\label{fig:slow-shock-fast-rarefaction}
\end{figure}

\begin{figure}[H]
\centering
\includegraphics[width=.78\linewidth]{figures/slow-rarefaction-fast-rarefaction-source.png}
\caption{Slow rarefaction--fast rarefaction.}\label{fig:slow-rarefaction-fast-rarefaction}
\end{figure}
